\documentclass[11pt]{article}
\usepackage{explorer}
\newcommand{\ExplorerLanguage}{FR}
\newcommand{\ExplorerNumber}{12}
\newcommand{\ExplorerKind}{Leçon et méthode}
\begin{document}
\ExplorerTitle{Terminale · analyse et algorithmique}{Encadrer une racine par dichotomie}
\ExplorerSection{Avant de commencer}
Continuité, signe et intervalles.\par
\ExplorerSection{Idées et méthode}
Si f est continue sur [a,b] et si f(a) et f(b) sont de signes opposés, le théorème des valeurs intermédiaires garantit au moins une racine dans [a,b]. Il ne garantit pas l’unicité.\par
On calcule \(m=\frac{a+b}{2}\). Si \(f(m)=0\), on a trouvé une racine. Sinon, on garde la moitié dont les valeurs aux bornes ont des signes opposés. La continuité maintient la garantie à chaque étape.\par
Après n partages, la largeur de l’intervalle est \(\frac{b-a}{2^{n}}\). Son milieu approche une racine avec une erreur au plus égale à la moitié de cette largeur. La largeur n’est donc pas la même chose que l’erreur du milieu.\par
\ExplorerSection{Exemple expliqué}
Pour \(f(x)=x^{2}-2\) sur [1,2], f est continue et strictement croissante. Sa racine est unique. Les milieux 1,5, puis 1,25 donnent les intervalles [1 ; 1,5], puis [1,25 ; 1,5] : on conserve toujours un signe négatif à gauche et positif à droite.\par
\ExplorerSection{À explorer}
Sur une droite graduée, représenter les intervalles successifs autour de \(\sqrt{2}\). Inscrire leur largeur à chaque étape.\par
\end{document}
